%%%%%
%
\section{Conversion du format \texorpdfstring{\texttt{int}}{int}
vers {\texorpdfstring{\texttt{float}}{float}} }

\itemstyle{Op\'erateur~:} \operfmt{CnvInt2Flt}
\index{CnvInt2Flt}

%%%%%OBJET%%%%
\itemstyle{Objet~:}
Transcodage d'un param\`etre du format \texttt{int\_16} vers le 
format \texttt{float}

%%%%%ARGUMENTS%%%%
\bigskip
%%%%%ENTREES%%%%
\begin{tabarg}
\multicolumn{3}{l}{\titretab {Entr\'ees}}\\
\hline
$X_i$ & Vecteur, \texttt{Int\_16} & Param�tre � convertir (unit� indiff�rente)\\
\end{tabarg}

\medskip
%%%%%SORTIES%%%%
\begin{tabarg}
\multicolumn{3}{l}{\titretab {Sorties}}\\
\hline
$X_f$ & Vecteur, \texttt{float} & $X_i$ converti en \texttt{float} (m�me unit� que $X_i$)\\
\end{tabarg}


%%%%%FORMULE%%%%
\itemstyle{Formule}

%%CnvInt2Flt
\[
X_f = X_i
\]



%%%%%COMMENTAIRE%%%%
\itemstyle{Commentaire~:}
Cette op�ration ne comporte pas de risque de d�bordement puisque
les flottants 32 bits selon la norme IEEE754 ont une mantisse
de 23 bits, dont la pr�cision relative est largement suffisante 
pour des entiers 16 bits.

%%%%%AUTEUR%%%%
\itemstyle{Auteur~:}
CNRM/GMEI/TRAMM

%%%%%REFERENCE%%%%
\itemstyle{R\'ef\'erences~:}
